\chapter{Vorwort}

\section*{Eidesstattliche Erklärung}
Hier sollte eine kluge Eidesstattliche Erklärung stehen, dass du alles allein gemacht hast et cetera et cetera.
\bigskip\underline{\hspace{0.2\textwidth}}\hspace{0.1\textwidth}\underline{\hspace{0.4\textwidth}}
\section*{Danksagung}
An erster Stelle dankt man meistens dem Betreuer. Je nach Sentimentalität kann dann auch weiteren Personen aus dem akademischen oder dem privaten Umfeld gedankt werden.
\section{Einführung}
(Grundlagen der Kryptographie, vgl. S.2ff.)
Der Wunsch, Informationen vor dem Zugriff unbefugter Personen zu schützen,
reicht weit in die Geschichte zurück. Bereits Julius Caesar soll ein
Verschlüsselungsverfahren verwendet haben, um militärische Nachrichten
vertraulich zu übermitteln. Bei der heute als Caesar-Chiffre bezeichneten
Methode werden die Buchstaben des Alphabets um eine festgelegte Anzahl von
Stellen verschoben. Nur wer die verwendete Verschiebung kennt, kann die
Nachricht unmittelbar entschlüsseln \cite[S.~14--15]{buell24}. Aus heutiger
Sicht bietet dieses Verfahren keine ausreichende Sicherheit, veranschaulicht
jedoch bereits ein grundlegendes Prinzip der Kryptographie: Eine Nachricht
wird mithilfe eines Schlüssels so verändert, dass ihr ursprünglicher Inhalt
für Unbefugte nicht mehr erkennbar ist.

Mit der zunehmenden Digitalisierung hat die sichere Übertragung von
Informationen erheblich an Bedeutung gewonnen. Beim Onlinebanking, bei der
E-Mail-Kommunikation und bei der Nutzung von Nachrichtendiensten werden
regelmäßig sensible Daten über öffentliche Netzwerke übertragen.
Kryptographische Verfahren sollen unter anderem gewährleisten, dass diese
Informationen nicht von unbefugten Dritten gelesen oder unbemerkt verändert
werden können.

Eine wichtige Rolle spielt dabei die asymmetrische Kryptographie, bei der
öffentliche und private Schlüssel verwendet werden. Die Sicherheit der
entsprechenden Verfahren beruht häufig auf mathematischen Problemen, für die
auf klassischen Computern keine effizienten Lösungsverfahren bekannt sind.
Zu diesen Problemen gehört das diskrete Logarithmusproblem. Wird es in der
Punktgruppe einer elliptischen Kurve über einem endlichen Körper betrachtet,
bildet es eine wesentliche Sicherheitsgrundlage der
Elliptische-Kurven-Kryptographie.

Die Entwicklung von Quantencomputern stellt diese Sicherheitsannahme
grundsätzlich infrage. Insbesondere ermöglicht der Shor-Algorithmus,
diskrete Logarithmusprobleme auf einem hinreichend leistungsfähigen
Quantencomputer effizient zu lösen. Daher ist zu untersuchen, welche
Konsequenzen sich aus der Entwicklung von Quantencomputern für die Sicherheit
kryptographischer Verfahren auf elliptischen Kurven ergeben.

\section{Zielsetzung und Forschungsfrage}

Ziel dieser Arbeit ist es, elliptische Kurven mathematisch einzuführen und
ihre Verwendung in der Kryptographie darzustellen. Dazu werden elliptische
Kurven über endlichen Körpern sowie das diskrete Logarithmusproblem in ihren
Punktgruppen untersucht. Anschließend werden die Grundlagen des
Quantencomputings und insbesondere der Shor-Algorithmus betrachtet, um dessen
Auswirkungen auf die Sicherheit der Elliptische-Kurven-Kryptographie zu
bewerten.

Daraus ergibt sich die folgende Forschungsfrage:

\begin{quote}
    \emph{Wie verändert der Shor-Algorithmus die Komplexität des diskreten
        Logarithmusproblems auf elliptischen Kurven, und welche Konsequenzen
        ergeben sich daraus für die Sicherheit der
        Elliptische-Kurven-Kryptographie?}
\end{quote}

\section{Aufbau der Arbeit}

Zunächst werden die für die weitere Untersuchung benötigten Grundlagen der
asymmetrischen Kryptographie eingeführt. Das RSA-Kryptosystem dient dabei als
Beispiel, um das Zusammenspiel öffentlicher und privater Schlüssel sowie das
Prinzip der asymmetrischen Verschlüsselung zu erläutern. Anschließend wird das
diskrete Logarithmusproblem zunächst in einer allgemeinen Gruppe formuliert.

Darauf folgt die mathematische Einführung elliptischer Kurven. Neben ihrer
Definition werden insbesondere die Gruppenstruktur, die Punktaddition und die
Skalarmultiplikation behandelt. Anschließend werden elliptische Kurven über
endlichen Körpern untersucht, da deren endliche Punktgruppen die Grundlage
kryptographischer Anwendungen bilden.

Auf diesen Grundlagen aufbauend wird das diskrete Logarithmusproblem auf
elliptischen Kurven betrachtet und seine Bedeutung für die Sicherheit
kryptographischer Verfahren erläutert. Danach werden ausgewählte
kryptographische Verfahren auf elliptischen Kurven vorgestellt.

Im letzten Teil der Arbeit werden zunächst die erforderlichen Grundlagen des
Quantencomputings eingeführt. Anschließend wird der Shor-Algorithmus und seine
Anwendung auf das diskrete Logarithmusproblem untersucht. Abschließend werden
die daraus resultierenden theoretischen und praktischen Konsequenzen für die
Sicherheit der Elliptische-Kurven-Kryptographie bewertet.


\section{Mathematische Grundlagen und Notation}

Anmerkung: Es wird davon ausgegeganngen, dass der Leser grundlegende algebraische Strukturen kennt. Verwendete Strukturen sind: Gruppen, Ringe, Körper, Vektorräume 
\textbf{(siehe Quelle einfügen)}
Körper
Zyklische Gruppe und Generator 
Eulersche Phi Funktion mit Primfaktorzerlegung.
Homo- und Endomorphismen
Algebraischer Abschluss
Vektorraum 
\[
\begin{aligned}
\phi(n) &\quad \text{Euler-Funktion},\\
\mathbb{Z}_{\varphi(n)}^{*} &\quad \text{multiplikative Gruppe modulo } n.\\
	\simeq
\end{aligned}
\]